\scanpage{130}
\mspage{-III- \handwritten{5} -}
\begin{flushright}n°\handwritten{3\unsure{8|6}3}\end{flushright}
On a vu en effet que $C$ (ou une de ses transformées) est l'ensemble des $\displaystyle\handwritten{\sum}_{i\in I}a_i e'_i$ tels que $a_i>0$, et que $\cos(\handwritten{\widehat}{e'_i,e'_j})=\cos(\handwritten{\pi}/m_{ij})\handwritten{\geq}0$ (prop.2).

COROLLAIRE 2.- \uline{Pour tout} $H\in\underline{H}$, \uline{soit} $e_H$ \uline{le vecteur unitaire orthogonal à} $H$ \uline{et situé du même côté de} $H$ \uline{que} $C$. \uline{Alors} $e_H$ \uline{est combinaison linéaire à coefficients positifs de} $e_1,\ldots,e_m$.

En effet, d'après la prop.3, on a $\displaystyle e_H=\handwritten{\sum}_{j\in I}b_j e_j$ avec $b_j\in R$, d'où $b_j=(e_H\handwritten{|}e'_j)$. Comme $e'_j\in\overline{C}$, et que $e_H$ est du même côté de $H$ que $C$, on a $(e_H\handwritten{|}e'_j)\handwritten{\geq}0$.

COROLLAIRE 3.- \uline{Avec les mêmes notations}, $s_i$ \uline{laisse stable l'ensemble des} $e_H$ \uline{distincts de} $e_i$.

Pour $e_H\handwritten{\ne}e_i$, on a $\displaystyle e_H=\handwritten{\sum}_{j=1}^m a_j e_j$ où les $a_j$ sont des nombres réels positifs (cor.2), et où l'un des $a_j$, pour $j\handwritten{\ne}i$, est $>0$. Le vecteur $s_i(e_H)=e_H-2(e_H\handwritten{|}e_i)e_i$ est, soit un des $e_{H'}$, soit un de leurs opposés ; or, pour $j\handwritten{\ne}i$, le coefficient de $e_j$ dans la décomposition de $s_i(e_H)$ est $a_j$ ; donc $s_i(e_H)$ ne peut être de la forme $-e_{H'}$ car ces $a_j$ ne sont pas tous $\handwritten{\leq}0$ (cor.2). Ainsi $s_i(e_H)$ est de la forme $e_{H'}$. C.Q.F.D.

COROLLAIRE 4.- \uline{L'ordre de} $G$ \uline{est égal au nombre des chambres définies par} $\underline{H}$.

Ceci résulte aussitôt de b).

COROLLAIRE 5.- \uline{Avec les notations du th.1}, \uline{l'ensemble des réflexions} $s_H$ $(H\in\underline{H})$ \uline{est nécessairement l'ensemble de toutes les réflexions de} $G$.

Soit $\underline{H}'\handwritten{\supset}\underline{H}$ l'ensemble des hyperplans $H$ tels que $s_H\in G$. Toute chambre relative à $\underline{H}'$ est contenue dans une chambre relative à $\underline{H}$ et une seule. D'après le cor.4, les chambres relatives à $\underline{H}'$ et à $\underline{H}$ sont donc les mêmes. D'après le lemme 3 du \S\ 1, n°1, tout élément de $\underline{H}'$ est mur d'une chambre, donc appartient à $\underline{H}$.

COROLLAIRE 6.- \uline{Soient} $C$ \uline{et} $C'$ \uline{deux chambres}. \uline{Alors} $\overline{C}\handwritten{\cap}\overline{C'}$ \uline{est l'adhérence d'une face} $F$ \uline{de} $C$. \uline{L'élément} $g\in G$ \uline{tel que} $C'=g(C)$ \uline{est un produit de réflexions} $s_H$ $(H\in\underline{H})$ \uline{par rapport à des hyperplans} $H$ \uline{contenant} $F$.

C'est une traduction du lemme 7, \S\ 2, n°4.
% archive sheet number: 129
